\section{关键流程设计}

本章从系统运行角度说明用户提交资源后，平台如何完成资源登记、任务创建、异步检测、人工审核、报告输出与状态通知。与前几章的静态模块和数据设计不同，本章强调对象在时间维度上的流转顺序，以及不同模块之间的协作边界。

需要说明的是：当前代码主干已经完整实现图像上传、图像检测、人工审核和通知闭环；论文检测与 Review 检测尚未形成独立生产链路，因此本章对这两部分采用“迭代目标流程”的方式描述，并明确其与现有图像检测流程的衔接点。

\subsection{资源上传与检测流程}
资源上传与检测发起流程是整个平台的入口流程。当前实现中，系统支持用户上传单张图像、PDF 文档以及压缩包三类输入，其中 PDF 和压缩包会在上传阶段被解析为可检测的图像集合。随后用户从已登记资源中选择图像发起检测，系统完成任务创建、额度校验和异步调度。

\begin{figure}[H]
    \centering
    \resizebox{0.98\textwidth}{!}{\begin{tikzpicture}[
  node distance=8mm and 8mm,
  font=\small,
  proc/.style={draw, rounded corners=2pt, fill=blue!6, align=center, minimum width=2.8cm, minimum height=0.9cm},
  store/.style={draw, rounded corners=2pt, fill=yellow!10, align=center, minimum width=2.8cm, minimum height=0.9cm},
  ext/.style={draw, rounded corners=2pt, fill=green!8, align=center, minimum width=2.8cm, minimum height=0.9cm},
  flow/.style={-{Latex[length=2.2mm]}, thick}
]

\node[ext] (user) {用户端};
\node[proc, right=of user] (upload) {上传文件\\权限校验};
\node[store, right=of upload] (file) {登记文件\\FileManagement};
\node[proc, right=of file] (parse) {按类型预处理\\图像/PDF/压缩包};
\node[store, right=of parse] (image) {生成图像资源\\ImageUpload};

\node[proc, below=of file] (select) {选择图像并提交\\任务参数};
\node[proc, right=of select] (quota) {校验提交权限\\与组织配额};
\node[store, right=of quota] (task) {创建 DetectionTask\\DetectionResult};
\node[proc, right=of task] (batch) {批量打包 img.zip\\和 data.json};
\node[ext, right=of batch] (celery) {Celery 队列};

\node[ext, below=of batch] (ai) {AI 推理服务};
\node[proc, left=of ai] (finalize) {结果回写\\报告生成\\任务完成通知};

\draw[flow] (user) -- (upload);
\draw[flow] (upload) -- (file);
\draw[flow] (file) -- (parse);
\draw[flow] (parse) -- (image);
\draw[flow] (image) |- (select);
\draw[flow] (select) -- (quota);
\draw[flow] (quota) -- (task);
\draw[flow] (task) -- (batch);
\draw[flow] (batch) -- (celery);
\draw[flow] (celery) -- (ai);
\draw[flow] (ai) -- (finalize);
\draw[flow] (finalize) |- (user);

\end{tikzpicture}
}
    \caption{资源上传与检测发起流程图}
\end{figure}

\subsubsection{流程输入}
\begin{itemize}[leftmargin=2em]
    \item 用户登录态、用户角色和权限信息；
    \item 原始上传资源，包括图像文件、PDF 文件或压缩包；
    \item 用户选择的检测图片集合、任务名称、检测参数和可选模型/检测模式。
\end{itemize}

\subsubsection{处理步骤}
\begin{enumerate}[leftmargin=2em]
    \item 用户上传文件后，后端首先校验其上传权限，随后在 \texttt{FileManagement} 中登记文件元数据；
    \item 系统将原始文件写入媒体存储目录，并根据文件类型进入不同的预处理分支；
    \item 若上传对象为单张图像，则直接创建对应的 \texttt{ImageUpload} 记录；若为 PDF，则提取其中图像页并批量生成 \texttt{ImageUpload}；若为压缩包，则依次处理其中的图像和 PDF；
    \item 上传完成后，系统写入上传日志，前端通过文件详情接口或图像列表接口获取当前可检测图像集合；
    \item 用户在前端选择待检测图像并提交任务，后端校验提交权限、组织配额、检测参数和所选模型配置；
    \item 系统创建 \texttt{DetectionTask}，记录任务类型、检测模式、模型版本和参数快照，并为选中图像逐一建立 \texttt{DetectionResult} 初始记录，状态置为处理中；
    \item 后端按批次将图像打包为临时压缩文件和参数文件，并将批次任务提交给 Celery 队列等待异步执行。
\end{enumerate}

\subsubsection{流程输出}
\begin{itemize}[leftmargin=2em]
    \item 已登记的文件记录和图像记录；
    \item 处于待处理或处理中状态的检测任务；
    \item 后续异步检测阶段所需的批处理上下文。
\end{itemize}

\subsubsection{关键控制点}
\begin{itemize}[leftmargin=2em]
    \item 上传阶段必须先进行权限判定，避免无上传权限用户占用存储空间；
    \item 检测提交前需完成组织额度校验，并在校验通过后立即扣减，防止并发下超额提交；
    \item 模型配置读取应以任务创建时刻的配置快照为准，避免模型启停或默认策略变更影响运行中的任务；
    \item 资源预处理和任务提交之间通过数据库对象衔接，避免前端直接依赖临时文件路径。
\end{itemize}

\subsection{图像检测流程}
图像检测流程描述的是异步任务真正执行的主链路。当前项目采用“批量发送到 AI 服务，再按单图并行后处理，最后任务级汇总收尾”的三阶段模式，兼顾了远程推理服务的调用效率和结果落库的并发能力。

\begin{figure}[H]
    \centering
    \resizebox{0.98\textwidth}{!}{\begin{tikzpicture}[
  node distance=8mm and 8mm,
  font=\small,
  proc/.style={draw, rounded corners=2pt, fill=blue!6, align=center, minimum width=3.1cm, minimum height=0.9cm},
  store/.style={draw, rounded corners=2pt, fill=yellow!10, align=center, minimum width=3.1cm, minimum height=0.9cm},
  ext/.style={draw, rounded corners=2pt, fill=green!8, align=center, minimum width=3.1cm, minimum height=0.9cm},
  flow/.style={-{Latex[length=2.2mm]}, thick}
]

\node[store] (task) {检测任务\\DetectionTask};
\node[proc, right=of task] (batch) {批次任务\\fetch\_batch};
\node[ext, right=of batch] (ai) {AI 推理服务\\批量返回};
\node[proc, right=of ai] (single) {单图后处理\\process\_single\_result};
\node[store, right=of single] (result) {结果落库\\DetectionResult\\SubDetectionResult};

\node[proc, below=of single] (final) {任务收尾\\finalize\_task};
\node[store, left=of final] (report) {生成检测报告};
\node[ext, right=of final] (notify) {通知用户\\WebSocket};

\draw[flow] (task) -- (batch);
\draw[flow] (batch) -- (ai);
\draw[flow] (ai) -- (single);
\draw[flow] (single) -- (result);
\draw[flow] (result) |- (final);
\draw[flow] (final) -- (report);
\draw[flow] (final) -- (notify);

\end{tikzpicture}
}
    \caption{图像检测执行流程图}
\end{figure}

\subsubsection{流程阶段划分}
\begin{enumerate}[leftmargin=2em]
    \item \textbf{批量获取阶段：} Celery 的 \texttt{fetch\_batch} 任务读取批次目录中的 \texttt{img.zip} 和参数文件，向远程 AI 推理服务发起整批请求；
    \item \textbf{单图后处理阶段：} 对每个 \texttt{DetectionResult} 调用 \texttt{process\_single\_result}，完成可视化图片保存、主结果更新、子方法结果写入和图像状态更新；
    \item \textbf{任务收尾阶段：} 当同一检测任务下所有图片的结果均完成后，执行 \texttt{finalize\_task}，统一回写任务完成时间、生成检测报告并向用户推送完成通知。
\end{enumerate}

\subsubsection{详细处理步骤}
\begin{enumerate}[leftmargin=2em]
    \item 异步任务启动后，将 \texttt{DetectionTask} 状态更新为进行中，并将本批次对应的 \texttt{DetectionResult} 状态标记为处理中；
    \item 系统调用远程 AI 推理服务获取整批检测结果，必要时执行重连与重试；
    \item 对批次中的每张图，系统分别解析总体真假判定、置信度、LLM 文本说明、ELA 可视化、EXIF 判定和各子方法掩码结果；
    \item 后处理任务将可视化图像保存到媒体目录，并更新 \texttt{DetectionResult}、\texttt{SubDetectionResult} 与 \texttt{ImageUpload}；
    \item 系统统计同一任务下已完成结果数量，若全部完成，则将任务状态更新为已完成并生成任务级 PDF 报告；
    \item 任务完成后，通过 WebSocket 群组向用户发送检测完成通知，前端据此刷新任务状态或提示用户查看结果。
\end{enumerate}

\subsubsection{流程输出}
\begin{itemize}[leftmargin=2em]
    \item 图像级检测结论、可视化结果和子方法输出；
    \item 已完成状态的检测任务和可下载的任务级报告；
    \item 面向用户端的任务完成通知。
\end{itemize}

\subsubsection{异常与恢复}
\begin{itemize}[leftmargin=2em]
    \item 当 AI 推理服务不可达或返回数量异常时，批处理任务执行重试与重连逻辑；
    \item 若个别图片后处理失败，不应提前将整个任务置为完成状态；
    \item 若报告文件尚未生成，则报告下载接口返回“正在生成”而不是直接报错。
\end{itemize}

\subsection{论文检测流程}
论文检测属于本轮迭代的目标能力。当前代码尚未形成独立的论文解析、段落切分和论文级结果实体，但根据需求文档和现有图像链路，可以将论文检测流程设计为“论文资源登记 -> 文本/图像双通道解析 -> 段落级分析 -> 论文级结果聚合”的模式。

\begin{figure}[H]
    \centering
    \resizebox{0.98\textwidth}{!}{\begin{tikzpicture}[
  node distance=8mm and 8mm,
  font=\small,
  proc/.style={draw, rounded corners=2pt, fill=blue!6, align=center, minimum width=3.0cm, minimum height=0.9cm},
  store/.style={draw, rounded corners=2pt, fill=yellow!10, align=center, minimum width=3.0cm, minimum height=0.9cm},
  ext/.style={draw, rounded corners=2pt, fill=green!8, align=center, minimum width=3.0cm, minimum height=0.9cm},
  flow/.style={-{Latex[length=2.2mm]}, thick}
]

\node[ext] (upload) {论文上传\\PDF/文档};
\node[store, right=of upload] (resource) {论文资源登记};
\node[proc, right=of resource] (parse) {结构解析\\正文/章节/图片};
\node[proc, right=of parse] (split) {文本切分\\段落级输入};

\node[proc, below=of split] (text) {AIGC 倾向\\事实性分析};
\node[proc, left=of text] (image) {插图检测\\复用图像链路};
\node[store, left=of image] (result) {论文级结果聚合};
\node[ext, left=of result] (report) {论文检测报告};

\draw[flow] (upload) -- (resource);
\draw[flow] (resource) -- (parse);
\draw[flow] (parse) -- (split);
\draw[flow] (split) -- (text);
\draw[flow] (parse.south) |- (image.north);
\draw[flow] (text) -- (image);
\draw[flow] (image) -- (result);
\draw[flow] (result) -- (report);

\end{tikzpicture}
}
    \caption{论文检测目标流程图}
\end{figure}

\subsubsection{目标流程}
\begin{enumerate}[leftmargin=2em]
    \item 用户上传论文文件后，系统先完成文件登记和基础格式校验；
    \item 解析模块提取论文正文、章节结构、元信息和内嵌图像，并建立论文资源与图像资源的关联；
    \item 文本分析模块对论文进行段落切分，执行 AIGC 倾向分析、事实性说明生成和异常段落标记；
    \item 若论文中包含图像资源，则可选择复用现有图像检测链路，对论文插图执行独立鉴伪；
    \item 聚合模块将文本检测结果、插图检测结果和整体风险评分整合为论文级结论，并将可参与综合融合的结构化子结果传递给报告服务模块；
    \item 报告模块输出论文检测报告，并在用户端展示高风险段落、对应说明和插图结果摘要。
\end{enumerate}

\subsubsection{与当前实现的衔接点}
\begin{itemize}[leftmargin=2em]
    \item 资源登记仍可复用现有文件上传和媒体存储机制；
    \item 任务状态流转、异步调度、报告生成和通知机制可复用当前图像检测任务框架；
    \item 需要新增的主要能力集中在论文解析、段落级结果建模和论文报告聚合。
\end{itemize}

\subsubsection{设计关注点}
\begin{itemize}[leftmargin=2em]
    \item 论文检测输出不应只有单一真假结论，而应保留段落级风险与解释信息；
    \item 论文文本链路和论文插图链路需要统一汇总口径，避免生成两套割裂的报告；
    \item 若论文检测与图像检测共同参与同一任务，需要保留跨模态一致性校验所需的资源关系和结果映射；
    \item 论文检测属于迭代能力，本阶段流程设计应强调扩展接口和复用现有基础设施。
\end{itemize}

\subsection{多模态联合检测与综合融合流程}
多模态联合检测流程用于处理同一任务中同时存在论文文本、学术图像、Review 文本或人工审核结果的场景。该流程不替代各子检测模块，而是在子模块完成后进行统一汇总、跨模态一致性校验和综合风险解释。

\subsubsection{处理步骤}
\begin{enumerate}[leftmargin=2em]
    \item 检测任务管理模块收集同一任务下的图像检测结果、论文检测结果、Review 检测结果和人工审核结果；
    \item 报告服务模块按任务类型读取各子模块结构化输出，包括图像可疑区域、论文高风险段落、事实性子结论、Review 模板化片段和人工审核意见；
    \item 系统执行跨模态一致性校验，例如论文文本结论与插图风险、Review 内容与论文主题、人工审核意见与自动检测结论之间是否存在明显冲突；
    \item 融合策略根据各维度风险分值、证据数量、人工审核权重和模型配置生成综合可信度评分与风险等级；
    \item 报告服务模块输出综合鉴伪结论、分维度风险说明、可疑内容定位、冲突说明和人工复核建议；
    \item 用户端在综合报告或检测记录详情中展示融合结论，同时保留各子模块原始证据入口。
\end{enumerate}

\subsubsection{关键控制点}
\begin{itemize}[leftmargin=2em]
    \item 多模态融合必须保留各子检测结果的原始结论和模型版本，避免综合分数掩盖证据来源；
    \item 跨模态一致性校验输出应作为风险提示和解释信息，不应直接覆盖人工审核或子模型原始结果；
    \item 融合策略、权重和风险等级阈值应可配置，并在报告中体现必要的解释说明。
\end{itemize}

\subsection{Review检测流程}
Review 检测同样属于迭代目标能力。该流程面向期刊审稿意见、同行评审文本或内部评阅意见，目标是识别模板化倾向、AI 生成倾向和文本表达异常。

\begin{figure}[H]
    \centering
    \resizebox{0.98\textwidth}{!}{\begin{tikzpicture}[
  node distance=8mm and 8mm,
  font=\small,
  proc/.style={draw, rounded corners=2pt, fill=blue!6, align=center, minimum width=3.0cm, minimum height=0.9cm},
  store/.style={draw, rounded corners=2pt, fill=yellow!10, align=center, minimum width=3.0cm, minimum height=0.9cm},
  ext/.style={draw, rounded corners=2pt, fill=green!8, align=center, minimum width=3.0cm, minimum height=0.9cm},
  flow/.style={-{Latex[length=2.2mm]}, thick}
]

\node[ext] (input) {Review 文本\\或文件输入};
\node[proc, right=of input] (clean) {文本清洗\\分段与校验};
\node[proc, right=of clean] (template) {模板化倾向\\相似性分析};
\node[proc, right=of template] (ai) {AI 生成倾向\\文本特征分析};

\node[proc, below=of ai] (evidence) {证据提取\\可疑片段高亮};
\node[store, left=of evidence] (risk) {风险评分\\结论说明};
\node[ext, left=of risk] (output) {结果展示\\报告摘要};

\draw[flow] (input) -- (clean);
\draw[flow] (clean) -- (template);
\draw[flow] (template) -- (ai);
\draw[flow] (ai) -- (evidence);
\draw[flow] (evidence) -- (risk);
\draw[flow] (risk) -- (output);

\end{tikzpicture}
}
    \caption{Review 检测目标流程图}
\end{figure}

\subsubsection{目标流程}
\begin{enumerate}[leftmargin=2em]
    \item 用户提交 Review 文本或上传包含 Review 内容的文件；
    \item 系统完成文本清洗、编码统一、结构化分段和长度校验；
    \item 检测模块分别执行模板化相似性分析、AI 生成倾向分析、术语密度与意见覆盖度分析；
    \item 聚合模块根据多项指标生成 Review 风险评分与判定说明；
    \item 结果层向前端返回总体结论、风险证据和可疑片段高亮信息；
    \item 报告层生成可归档的 Review 检测结果摘要，供期刊管理或后续人工复核使用。
\end{enumerate}

\subsubsection{与当前实现的衔接点}
\begin{itemize}[leftmargin=2em]
    \item 任务创建、异步执行、通知和日志记录可沿用现有检测任务框架；
    \item 需要新增 Review 资源表示、文本分析器和 Review 结果展示组件；
    \item 若 Review 检测结论需要人工复核，可复用当前人工审核模块的审批与指派思想。
\end{itemize}

\subsubsection{设计关注点}
\begin{itemize}[leftmargin=2em]
    \item Review 检测更强调“风险提示”而非简单真假二分，需要保留证据型输出；
    \item 结果展示应适配文本高亮和说明片段，不宜直接套用图像检测结果页；
    \item 后续若引入模型管理，应记录 Review 检测所用模型版本与规则版本。
\end{itemize}

\subsection{人工审核流程}
人工审核流程用于处理用户对自动检测结果存在异议、或者系统要求进一步专家复核的场景。当前代码已经形成“发布者申请 -> 管理员审批 -> 审核员执行 -> 审核结果回传”的完整链路；迭代需求进一步要求支持公开审核任务池、评论点赞、举报处置和发起人取消审核。

\begin{figure}[H]
    \centering
    \resizebox{0.98\textwidth}{!}{\begin{tikzpicture}[
  node distance=8mm and 10mm,
  font=\small,
  proc/.style={draw, rounded corners=2pt, fill=red!8, align=center, minimum width=3.0cm, minimum height=0.9cm},
  store/.style={draw, rounded corners=2pt, fill=yellow!10, align=center, minimum width=3.0cm, minimum height=0.9cm},
  ext/.style={draw, rounded corners=2pt, fill=green!8, align=center, minimum width=2.6cm, minimum height=0.9cm},
  flow/.style={-{Latex[length=2.2mm]}, thick}
]

\node[ext] (publisher) {发布者};
\node[proc, right=of publisher] (request) {提交审核申请\\选择图像与审核员};
\node[store, right=of request] (rr) {创建 ReviewRequest};
\node[ext, right=of rr] (admin) {组织管理员};

\node[proc, below=of rr] (approve) {审批通过/拒绝};
\node[store, left=of approve] (manual) {创建 ManualReview\\ImageReview};
\node[ext, left=of manual] (reviewer) {审核员};

\node[proc, below=of manual] (review) {逐图评分、标注\\提交人工结论};
\node[store, right=of review] (aggregate) {聚合审核结果\\生成人工审核报告};
\node[ext, right=of aggregate] (result) {发布者查看结果};

\draw[flow] (publisher) -- (request);
\draw[flow] (request) -- (rr);
\draw[flow] (rr) -- (admin);
\draw[flow] (admin) -- (approve);
\draw[flow] (approve) -- node[above]{通过} (manual);
\draw[flow] (manual) -- (reviewer);
\draw[flow] (reviewer) -- (review);
\draw[flow] (review) -- (aggregate);
\draw[flow] (aggregate) -- (result);

\end{tikzpicture}
}
    \caption{人工审核流程图}
\end{figure}

\subsubsection{处理步骤}
\begin{enumerate}[leftmargin=2em]
    \item 发布者在检测结果页选择若干图片并提交人工审核申请，系统创建 \texttt{ReviewRequest}，记录申请理由、关联图片和指定审核员；
    \item 系统向组织管理员发送待审批提醒，当前实现中该提醒通过邮件发送；
    \item 管理员审核申请后，若拒绝，则更新审核请求状态并结束流程；若通过，则根据任务配置进入公开审核任务池或为指定审核员创建 \texttt{ManualReview}，并为每张图初始化 \texttt{ImageReview}；
    \item 具备审核资格的用户可在公开任务池按学科、任务类型、紧急程度和风险等级筛选任务，系统校验组织边界、权限和利益冲突后允许进入详情或领取任务；
    \item 系统向被分派或领取任务的审核员发送站内通知，审核员在前端查看待办任务并逐图、逐段或逐 Review 填写评分、理由、标注点和人工结论；
    \item 审核任务详情页允许用户发表评论或点赞，评论写入前执行敏感词过滤；用户发现违规任务或评论时可提交举报，举报进入管理端处理队列；
    \item 发起人可在任务未完成前提交取消/终止请求，系统更新审核请求状态、阻止新的审核提交，并通知已参与审核者；
    \item 审核员提交后，系统保存逐图审核数据并更新人工审核状态；
    \item 当发布者或管理端查看审核结果时，系统聚合所有已完成审核记录，展示不同审核员的判定结果，并在需要时生成人工审核报告。
\end{enumerate}

\subsubsection{流程输出}
\begin{itemize}[leftmargin=2em]
    \item 审核申请记录、审核执行记录和逐图审核结果；
    \item 公开任务池列表、评论点赞记录、举报记录和取消审核记录；
    \item 审核员侧待办通知和发布者侧审核完成结果；
    \item 可下载的人工审核 PDF 报告。
\end{itemize}

\subsubsection{关键控制点}
\begin{itemize}[leftmargin=2em]
    \item 审核申请仅允许具备发布权限且角色为发布者的用户发起；
    \item 审核员范围需要受组织、角色、任务可见性和利益冲突规则约束，避免跨组织或不合规参与；
    \item 管理员审批和审核员执行属于两个独立状态层，流程设计中必须分开建模；
    \item 审核记录应保留逐图评分与理由，保证人工结论具备复核依据；
    \item 评论点赞和举报属于可审计交互行为，需要记录操作者、目标对象、时间和处理结果；
    \item 审核取消只能作用于未完成任务，取消后应保留既有审核意见和日志，避免历史证据丢失。
\end{itemize}

\subsection{通知与状态同步流程}
通知与状态同步流程用于保证用户在检测和审核过程中能够及时获取任务变化。当前实现采用“数据库持久化通知 + WebSocket 实时推送 + REST 查询补偿”的组合方案。

\begin{figure}[H]
    \centering
    \resizebox{0.98\textwidth}{!}{\begin{tikzpicture}[
  node distance=8mm and 9mm,
  font=\small,
  proc/.style={draw, rounded corners=2pt, fill=blue!8, align=center, minimum width=3.0cm, minimum height=0.9cm},
  store/.style={draw, rounded corners=2pt, fill=yellow!10, align=center, minimum width=3.0cm, minimum height=0.9cm},
  ext/.style={draw, rounded corners=2pt, fill=green!8, align=center, minimum width=2.7cm, minimum height=0.9cm},
  flow/.style={-{Latex[length=2.2mm]}, thick}
]

\node[proc] (event) {业务事件\\任务完成/审核分派/公告发布};
\node[store, right=of event] (db) {写入 Notification};
\node[proc, right=of db] (ws) {推送到用户通知组\\WebSocket};
\node[ext, right=of ws] (online) {在线前端};

\node[proc, below=of db] (rest) {通知查询接口\\未读数/列表/已读};
\node[ext, left=of rest] (offline) {离线或重连前端};

\draw[flow] (event) -- (db);
\draw[flow] (db) -- (ws);
\draw[flow] (ws) -- (online);
\draw[flow] (db) -- (rest);
\draw[flow] (offline) -- (rest);
\draw[flow] (rest) -- (db);

\end{tikzpicture}
}
    \caption{通知与状态同步流程图}
\end{figure}

\subsubsection{处理步骤}
\begin{enumerate}[leftmargin=2em]
    \item 业务模块在任务完成、人工审核分派、管理员审批、评论回复、点赞互动、举报处理、模型配置变更或公告发送时创建 \texttt{Notification} 记录；
    \item 在线用户通过带令牌参数的 WebSocket 连接加入对应的用户通知组；
    \item 当检测任务收尾阶段完成后，系统向用户对应的通知组推送实时消息，前端接收后刷新任务状态提示；
    \item 若用户不在线，则通知记录保留在数据库中，用户下次进入系统后可通过通知接口或社区反馈入口拉取未读消息；
    \item 用户可通过接口将全部通知或单条通知标记为已读，管理端则可以通过广播接口向所有普通用户发布系统公告。
\end{enumerate}

\subsubsection{流程特点}
\begin{itemize}[leftmargin=2em]
    \item WebSocket 负责降低用户等待期间的轮询压力，但系统不依赖 WebSocket 作为唯一通知通路；
    \item 通知记录落库后，即使实时推送失败，用户仍可通过接口补拉消息；
    \item 社区反馈入口是通知消息的业务聚合视图，需要按反馈类型关联到检测记录、审核任务、举报结果或报告详情；
    \item 当前实现的检测完成通知主要面向用户级群组，而不是面向单任务群组统一广播。
\end{itemize}

\subsection{流程图说明}
本章流程图采用“业务主线优先、异常分支从简”的表达策略。图中保留了对概要设计最关键的角色、模块和状态转换，不展开到接口参数或数据库字段级别。各图与正文的对应关系如下：
\begin{itemize}[leftmargin=2em]
    \item 图 7-1 对应资源上传与检测发起流程，体现资源登记、任务创建和异步调度的衔接；
    \item 图 7-2 对应图像检测执行流程，体现批量推理、单图后处理和任务收尾的异步链路；
    \item 图 7-3 对应论文检测目标流程，体现论文解析、文本检测、图像检测复用和论文级聚合；
    \item 图 7-4 对应 Review 检测目标流程，体现文本预处理、多维分析和风险解释；
    \item 多模态联合检测与综合融合流程体现图像、论文、Review 和人工审核结果进入综合报告前的聚合关系；
    \item 图 7-5 对应人工审核流程，体现发布者、管理员、审核任务池和审核员之间的协作关系；
    \item 图 7-6 对应通知与状态同步流程，体现数据库通知、社区反馈、实时推送和接口补偿三者的协同。
\end{itemize}

后续进入详细设计阶段时，可在本章基础上继续细化为接口时序图、状态机图和异常处理分支图，从而为实现与测试提供更直接的依据。
